\documentclass[a4paper,12pt]{article}
\usepackage{pdfpages}
\usepackage{ucs}
\usepackage[utf8x]{inputenc}
\usepackage[T1,T2A]{fontenc}
\usepackage{float}
\usepackage[russian]{babel}
\usepackage{graphicx}
\usepackage{listings}
\usepackage{xcolor}
\usepackage[left=20mm, top=15mm, right=15mm,bottom=15mm,nohead,footskip=10mm]{geometry}
\usepackage{indentfirst}
\renewcommand{\baselinestretch}{1}

\begin{document}
\thispagestyle{empty}
\begin{center}

\renewcommand{\baselinestretch}{1}
{\large
{\sc Министерство образования и науки Российской Федерации\\
ФГБОУ ВО <<Петрозаводский государственный университет>>\\
Институт математики и информационных технологий\\
Кафедра информатики и математического обеспечения
}
}

\end{center}

\vfill

\begin{center}
{\normalsize Отчет по учебной практике\\(компьютерные технологии в математике) } \\

%\bigskip
%	{\Large \sc Анализ производительности беспроводной сети 802.11} \\
\end{center}
\bigskip

\begin{flushright}

  \parbox{8cm} {

    \renewcommand{\baselinestretch}{1}

    \normalsize

    Выполнил:\\Петрушин А.~А. группа\ 22103 \\
%\begin{flushright}
%\vspace{5mm}
    \rule{55mm}{.3pt}\\
\mbox{\ \ \ \qqu}{\it подпись}\\
%\end{flushright}

Руководитель практики:\\
Е. И. Рыбин\\
%Оценка руководителя:\\ \hfill \rule{35mm}{.3pt}\\
%\begin{flushright}
%\vspace{5mm}
\rule{55mm}{.3pt}\\
\mbox{\ \ \ \qqu}{\it подпись}\\
%\end{flushright}

Итоговая оценка:\\
%\begin{flushright}
%<<\rule{10mm}{.3pt}>> \rule{35mm}{.3pt} 2022 г.\\
%\end{flushright}
%\begin{flushright}
%\vspace{5mm}
\rule{55mm}{.3pt}\\
\mbox{\ \ \ \qqu}{\it оценка}\hfill\\
%\end{flushright}

}
 
\end{flushright}

\vfill

\begin{center}
\large Петрозаводск -- 2024
\end{center}
\newpage
\tableofcontents
\newpage
\section{Описание работы}
 Первое задание учебной практики заключалось в том, что мы должны транслировать файл <<first-example.tex>> и получить документ в форматах dvi, postscript и pdf. Чтобы проделать данную работу мы ввели несколько команд в терминале, а именно: <<pdflatex first-example.tex>> для получения pdf-файла, <<latex first-example.tex>> для dvi-файла и <<dvips first-example.dvi>> для postscript-файла.

Вторая задание требовало подготовить документ, содержащий математический текст, который мы берём из 1 или 2 тома <<Краткого курса математического анализа>> Л. Д. Кудрявцева. Мне были предоставлены страницы с 178 до 180 из второго тома. Пользуясь пособиями, предоставленными на сайте кафедры Информатики и Математического Обеспечения ПетрГУ, я изучил синтаксис языка latex и сделал документ, который содержит предоставленный текст.

Последнее третье задание запрашивало построить изображение какой-либо функции в декартовых координатах с помощью интерпретатора gnuplot и разместить его в документе задания 2. Я ознакомился со скринкастами на сайте и сделал изображение графика функции:
$$x^2-y^2 $$
Данное изображение было сделано с помощью ряда таких команд и вставлено в конец документа:

1) set term pdfcairo

2) set output 'surface.pdf'

3) set ztics 50

4) set cbtics 50

5) splot x**2-y**2 t'Гиперболический парабалоид'w pm3d

\newpage
\section{Результаты работы}
\subsection{Задание 2}
\begin{figure}[H]
\centering
\includegraphics[width=1\textwidth]{2lab.png}
\caption{Код/Результат второго задания}
\end{figure}
\subsection{Задание 3}
\begin{figure}[h]
    \centering
    \includegraphics[width=0.8\textwidth]{surface.png}
    \caption{График функции $x^2 - y^2$}
\end{figure}

\end{document}
